\section{问题四：考虑尺寸变化的烘干过程}

\subsection{模型的建立}

\steptitle{尺寸收缩的实测过程}

实际烘干中水分流失会引起药材收缩。附件 2 给出半径由 $2.000$~cm 单调收缩到
\qfourRatdryingendcm 附近的完整过程，总收缩量为 \qfourshrinkagepercent；
前 $24$~h 完成全部收缩量的 $99.4\%$，此后半径基本不变。
物性改按附录 4 取为 $\rho=760+90C$、$c_p=1850+2150C/(C+1)$、$k=0.12+0.20C/(C+1)$、
$D=4.2\times10^{-4}e^{-0.30/C}e^{-3850/T}$（式中温度按附录原式取开尔文），
判据与问题三相同，仍为全域最大干基含水率降至 $0.15$~kg/kg 以下。

\steptitle{归一化坐标与表观对流项}

与前三问的根本差别在于计算域本身随时间内移。若仍在固定网格上求解，
边界处的通量将失去守恒性，收缩会在边界附近引入人工质量源。
因此我们引入\textbf{归一化坐标} $\xi=r/R(t)$，把移动边界化为固定域 $\xi\in[0,1]$，
按第五章式~\eqref{eq:moving} 得到

\begin{equation}
\left.\frac{\partial C}{\partial t}\right|_{\xi}
=\frac{1}{\xi R^{2}}\frac{\partial}{\partial \xi}
\left(D\,\xi\,\frac{\partial C}{\partial \xi}\right)
+\underbrace{\frac{\xi\dot R}{R}\frac{\partial C}{\partial \xi}}_{\text{表观对流项}}
\label{eq:q4moving}
\end{equation}

能量方程与之同形，系数换成热扩散率 $\alpha=k/(\rho c_p)$。
表观对流项完全已知：$\dot R$ 由附件 2 的单调插值解析求导得到
（改用线性插值只使烘干时长变化 \qfourradiusinterphours）。
\steptitle{报告层的取值方式}

报告层则按\textbf{当前物理距离}取值：每个时刻按 $r=\xi R(t)$ 映射回物理坐标，
$r>R(t)$ 的位置留空；末列“药材表面”取 $\xi=1$ 处的原生解变量。

\subsection{结果与分析}

\steptitle{烘干时刻与读数}

\textbf{全域最大干基含水率于 \qfourdryingtimehours 降至判据以下，此时半径已收缩到 \qfourRatdryingendcm。}
表~\ref{tab:table6_q4} 给出题目要求的表~6。$6$~h 时中心与表面的干基含水率分别为
\qfourtablesixtsixcenter 与 \qfourtablesixtsixsurface；烘干结束时分别为 \qfourtablesixendcenter
与 \qfourtablesixendsurface。

\begin{table}[H]
  \centering
  \caption{考虑尺寸变化时烘干过程的干基水分浓度（kg/kg）}
  \label{tab:table6_q4}
  \small
  \begin{tabular}{lccccc}
    \toprule
    时间/h & $r=0$ & $0.5$~cm & $1.0$~cm & $1.5$~cm & 药材表面 \\
    \midrule
    6 & 1.9473 & 1.7782 & 1.2540 & -- & 0.5654 \\
    12 & 0.8718 & 0.7750 & 0.4851 & -- & 0.2140 \\
    18 & 0.4588 & 0.4126 & 0.2650 & -- & 0.1019 \\
    24 & 0.3070 & 0.2803 & 0.1901 & -- & 0.0713 \\
    30 & 0.2379 & 0.2197 & 0.1554 & -- & 0.0610 \\
    36 & 0.1999 & 0.1860 & 0.1356 & -- & 0.0567 \\
    42 & 0.1760 & 0.1647 & 0.1226 & -- & 0.0545 \\
    48 & 0.1596 & 0.1499 & 0.1135 & -- & 0.0532 \\
    \textbf{52.6361} & \textbf{0.1500} & \textbf{0.1413} & \textbf{0.1081} & \textbf{--} & \textbf{0.0526} \\
    \bottomrule
  \end{tabular}
  \par\smallskip\footnotesize “--”表示该时刻 $r>R(t)$，该位置在药材之外。
\end{table}


\steptitle{收缩对报告列的影响}

值得注意的是，$1.5$~cm 一列在 $6$~h 时已经为空——该时刻半径已小于 $1.5$~cm，
按当前物理坐标的取法，该位置已经不在药材内，故按题目要求留空。
这说明尺寸收缩不仅改变扩散路径长度，还会改变“同一物理位置”在不同时刻是否属于药材，
在按固定半径列表时必须逐时刻判断。

图~\ref{fig:q4} 给出收缩曲线、与不收缩情形的对比，以及表面含水率的对比。

\begin{figure}[H]
  \centering
  \includegraphics[width=0.98\textwidth]{fig04_shrinkage_effect.pdf}
  \caption{尺寸收缩对烘干时长的影响}
  \label{fig:q4}
\end{figure}


\subsection{尺寸收缩对烘干时长的影响}

\steptitle{收缩为何缩短时长}

收缩与换物性这两个效应方向相反，量级对比才是关键。
附录 4 的扩散系数前因子（$4.2\times10^{-4}$）比附录 3（$2.4\times10^{-3}$）小 $5.7$ 倍，
仅换物性应当\textbf{拉长}时长；而收缩把扩散路径缩短到 $R_{\min}/R_0\approx0.599$，
按特征时间正比于半径平方，扩散加快到 $1/0.599^2\approx2.8$ 倍，压倒了前者的拖慢。
两者相抵的结果是 \qfourdryingtimehours，比问题三的 \qthreedryingtime 缩短
\qfourvsqthreedryingtimehours。\textbf{故收缩不可忽略，问题四的时长不能由问题三推算}；
沿用固定半径模型会把时长高估近五小时，这正是移动边界模型的必要性所在。
完整时空分布见图~\ref{fig:q4map}。

\begin{figure}[H]
  \centering
  \includegraphics[width=0.95\textwidth]{figA3_q4_shrinking_field.pdf}
  \caption{问题四收缩域含水率场（附录图）}
  \label{fig:q4map}
\end{figure}

\subsection{数值检验与适用性分析}

\steptitle{守恒恒等式与解析不变量检验}

移动边界下的守恒性难以目测，故按第五章式~\eqref{eq:conservation} 的离散恒等式逐时步核验：
$\frac{d}{dt}\sum_j \omega_j C_j
=-2\frac{\dot R}{R}\sum_j \omega_j C_j+\frac{\dot R}{R}C_N-\frac{h_m}{R}(C_N-C_\infty)$，
实测残差 \qfourmassbalancerel。
更锋利的判据是\textbf{解析不变量检验}：关闭表面传质、只让计算域收缩，干基含水率必须严格不变，
实测最大偏离仅 \qfourmovingboundarydev。表观对流项的符号或系数一旦写错，该检验会立刻暴露，
故它直接证明了式~\eqref{eq:q4moving} 新增项的正确性。

\steptitle{数值层面的其余证据}

其余三项与前三问同法：与 Bessel 级数解析解对照（固定半径、常数扩散系数的极限），
最大绝对误差 \qfourreferencemaxabserr；网格由 $400$ 加密到 $800$，
报告点最大差 \qfourgridrefinementdiff；内部步长减半，烘干时长变化 \qfourtimerefinementhours。
独立复算方面，格心实现（表面值由 Robin 半单元关系回算、不含启动步）与本文方法在表~6 上
最大差 \qfourbaselinetablesixdiff、时长相差 \qfourbaselinedryingtimeseconds，
差异集中在早期表面列，是其一阶回算表面值的代价。全部检验量汇总于表~\ref{tab:verify}，
图形汇总见图~\ref{fig:verify}。

\begin{table}[H]
  \centering
  \caption{四问的检验量汇总}\label{tab:verify}
  \small
  \setlength{\tabcolsep}{4pt}
  \begin{tabular*}{\textwidth}{@{\extracolsep{\fill}}l *{5}{c}@{}}
    \toprule
    问题 & 解析对照/(kg/kg) & 网格加密/(kg/kg) & 时间加密/h & 独立实现差/s & 离散守恒 \\
    \midrule
    第 (1) 问 & $6.18\times10^{-7}$ & $3.24\times10^{-6}$ & $3.89\times10^{-10}$ & $9.72\times10^{-7}$ & $5.41\times10^{-11}$ \\
    第 (2) 问 & $5.98\times10^{-7}$ & $3.2\times10^{-7}$   & $2.78\times10^{-4}$  & $26$              & $4.37\times10^{-13}$ \\
    第 (3) 问 & $5.98\times10^{-7}$ & $3.2\times10^{-7}$   & $3.6\times10^{-3}$   & $26$              & $4.76\times10^{-13}$ \\
    第 (4) 问 & $1.15\times10^{-6}$ & $2.12\times10^{-6}$  & $0$                  & $15$              & $2.56\times10^{-14}$ \\
    \bottomrule
  \end{tabular*}
  \par\smallskip\footnotesize 注：第 (3) 问的解析对照取与其同模型的第 (2) 问结果；
  第 (1) 问的时间加密与独立实现差两项为 kg/kg；离散守恒为无量纲相对残差。
\end{table}

\begin{figure}[H]
  \centering
  \includegraphics[width=0.92\textwidth]{fig05_numerical_verification.pdf}
  \caption{数值可信度的四类证据}
  \label{fig:verify}
\end{figure}

\steptitle{本问结论的适用条件}

本问以题目给定的 $h_m$、$h$、附录 4 经验式与“收缩严格按附件 2 进行”为条件，
质量方程取 Fick 标准形式（几何收缩在模型内不是质量汇），
口径与轴对称简化沿用问题二、问题三的说明，不再重复。
另有两点需专为声明：改用湿基自洽式会使烘干时长变化 \qfoursensmassequationformhours，
量级超过任何一项参数扰动，可见\textbf{本问的主导不确定性来自模型设定而非题目给定的传质系数}，
把不确定性归因于“参数不准”并不准确；烘干时长精确落在判据穿越时刻，
继承了判据尾部的平坦性。单项扰动中，$h_m$ 放大 $5\%$ 使时长变化 \qfoursenshmplusfivepcthours，
$D$ 放大 $5\%$ 变化 \qfoursensDplusfivepcthours，$T_\infty$ 升高 $2$~$^\circ$C
变化 \qfoursensTinfplustwoChours，可见工艺参数的调节余地有限。